\documentclass[11pt,a4paper]{article}
\usepackage[margin=2cm]{geometry}
\usepackage{xeCJK}
\usepackage{fontspec}
\usepackage{graphicx}
\usepackage{booktabs}
\usepackage{array}
\usepackage{longtable}
\usepackage{multirow}
\usepackage{xcolor}
\usepackage{hyperref}
\usepackage{amsmath}
\usepackage{enumitem}
\usepackage{fancyhdr}
\usepackage{titlesec}

\hypersetup{
    colorlinks=true,
    linkcolor=blue!80!black,
    citecolor=blue!80!black,
    urlcolor=blue!80!black,
    pdfborder={0 0 0}
}

% 中文字體設置
\setCJKmainfont{Noto Serif CJK TC}[
  UprightFont = * Light,
  BoldFont = * Bold,
  ItalicFont = * Light,
  BoldItalicFont = * Bold
]
\setCJKsansfont{Noto Sans CJK TC}[
  UprightFont = * Light,
  BoldFont = * Bold
]
\setCJKmonofont{Noto Sans Mono CJK TC}
\setmainfont{Times New Roman}

% 頁眉設置
\pagestyle{fancy}
\fancyhf{}
\fancyhead[L]{重振床邊臨床診察\\Strategies to Reinvigorate Bedside Clinical Encounter}
\fancyhead[R]{\thepage}
\renewcommand{\headrulewidth}{0.4pt}

% 標題格式
\titleformat{\section}{\Large\bfseries\color{blue!80!black}}{\thesection}{1em}{}
\titleformat{\subsection}{\large\bfseries\color{blue!60!black}}{\thesubsection}{1em}{}

% 自定義顏色
\definecolor{headercolor}{RGB}{70,130,180}
\definecolor{subheadercolor}{RGB}{100,149,237}

\title{\textbf{\Large Strategies to Reinvigorate the Bedside\\
Clinical Encounter\\
重振床邊臨床診察的策略}}
\author{謝慕揚醫師 MD, PhD, FESC\\
根據 Brian T. Garibaldi \& Stephen W. Russell 整理}
\date{2025年11月}

\begin{document}

\maketitle

\section{概述}

本文為《新英格蘭醫學雜誌》醫學教育系列文章，探討現代醫療環境中床邊臨床技能衰退的問題，並提供六大策略協助臨床教育者重振床邊醫學文化。研究指出，當代醫學訓練者僅花費13\%的時間與病人直接接觸，此趨勢導致基本床邊技能下降、診斷錯誤增加、臨床結果惡化，並增加醫療成本。

\section{研究資訊}

\subsection{基本資料}
\textbf{標題：}重振床邊臨床診察的策略\\
\textbf{作者：}Brian T. Garibaldi, Stephen W. Russell\\
\textbf{機構：}Northwestern University Feinberg School of Medicine; University of Alabama at Birmingham\\
\textbf{期刊：}\textit{New England Journal of Medicine} (2025)\\
\textbf{DOI：}\href{https://doi.org/10.1056/NEJMra2500226}{10.1056/NEJMra2500226}\\
\textbf{發表日期：}2025年11月12日

\subsection{問題背景}

\begin{longtable}{p{4cm}p{11cm}}
\toprule
\textbf{現況問題} & \textbf{影響與後果} \\
\midrule
\endfirsthead
\textbf{現況問題} & \textbf{影響與後果} \\
\midrule
\endhead
\bottomrule
\endfoot
\endlastfoot

直接接觸時間減少 & 
\begin{minipage}[t]{11cm}
- 醫學生僅13\%時間與病人直接接觸\\
- 基本床邊持續下降\\
- 診斷能力受到影響
\end{minipage} \\
\midrule

診斷錯誤增加 & 
\begin{minipage}[t]{11cm}
- 超過50\%門診診斷錯誤源於病史詢問不良\\
- 理學檢查執行錯誤\\
- 過度依賴科技檢查
\end{minipage} \\
\midrule

醫療成本上升 & 
\begin{minipage}[t]{11cm}
- 過度依賴科技導致過度檢查\\
- 臨床技能衰退引發不必要調查\\
- 醫療資源浪費
\end{minipage} \\
\midrule

醫病關係弱化 & 
\begin{minipage}[t]{11cm}
- 醫學生與住院醫師同理心下降\\
- 執業醫師壓力與職業倦怠增加\\
- 醫病關係信任感降低
\end{minipage} \\
\midrule

健康不平等惡化 & 
\begin{minipage}[t]{11cm}
- 床邊時間不足影響弱勢族群\\
- 加劇醫療照護差異\\
- 邊緣化群體受到不成比例影響
\end{minipage} \\

\end{longtable}

\section{歷史演變：如何走到今天？}

\subsection{早期美國醫學教育}

\begin{itemize}
    \item \textbf{19世紀初：}個別師徒制（preceptorship）
    \item \textbf{19世紀中：}學位授予的講授式醫學院
    \item \textbf{20世紀初：}William Osler於Johns Hopkins Hospital建立床邊教學典範
    \item \textbf{Osler教學法：}「學生檢查病人、做出診斷、聽到病變肺部的囉音、感受腫瘤異常堅硬的質地」
\end{itemize}

\subsection{現代醫學教育的轉變}

\begin{longtable}{p{5cm}p{9cm}}
\caption{導致床邊教學衰退的因素} \\
\toprule
\textbf{因素} & \textbf{影響} \\
\midrule
\endfirsthead
\textbf{因素} & \textbf{影響} \\
\midrule
\endhead
\bottomrule
\endfoot
\endlastfoot

科技為基礎的檢查 & 
\begin{minipage}[t]{9cm}
- 診斷過程移向實驗室與放射科\\
- 造成醫師所見、所感、所聞不再準確的錯誤印象\\
- 實體診察價值被低估
\end{minipage} \\
\midrule

電子病歷系統(EHR) & 
\begin{minipage}[t]{9cm}
- 工作流程迫使醫師花更多時間面對「iPatient」\\
- 實際與病人相處時間減少\\
- 臨床資料數位化優先於實體接觸
\end{minipage} \\
\midrule

值班時數限制 & 
\begin{minipage}[t]{9cm}
- 經濟壓力下優先考量效率\\
- 病人處理量優先於床邊評估\\
- 教學時間被壓縮
\end{minipage} \\
\midrule

時間碎片化 & 
\begin{minipage}[t]{9cm}
- 與病人相處的有限時間日益分散\\
- 難以進行完整評估\\
- 連續性照護受到挑戰
\end{minipage} \\
\midrule

教學能力下降 & 
\begin{minipage}[t]{9cm}
- 具床邊教學能力的教師減少\\
- 形成惡性循環\\
- 技能傳承中斷
\end{minipage} \\

\end{longtable}

\section{重振床邊醫學文化的六大策略}

\subsection{策略概覽}

\begin{longtable}{p{5cm}p{9cm}}
\caption{六大核心策略} \\
\toprule
\textbf{策略} & \textbf{核心理念} \\
\midrule
\endfirsthead
\textbf{策略} & \textbf{核心理念} \\
\midrule
\endhead
\bottomrule
\endfoot
\endlastfoot

\textbf{1. 前往床邊觀察} & 
\begin{minipage}[t]{9cm}
觀察病人與學習者\\
病人在哪裡，診斷線索就在哪裡
\end{minipage} \\
\midrule

\textbf{2. 實證理學檢查} & 
\begin{minipage}[t]{9cm}
教導與實踐基於證據的理學檢查方法\\
假說驅動的診察流程
\end{minipage} \\
\midrule

\textbf{3. 刻意練習機會} & 
\begin{minipage}[t]{9cm}
創造有目的的技能練習機會\\
晨間迴診與專門教學時段
\end{minipage} \\
\midrule

\textbf{4. 科技輔助教學} & 
\begin{minipage}[t]{9cm}
運用科技強化臨床診察技能\\
數位聽診器、床邊超音波、人工智慧
\end{minipage} \\
\midrule

\textbf{5. 回饋與評估} & 
\begin{minipage}[t]{9cm}
尋求並提供臨床技能回饋\\
直接觀察與情境回饋
\end{minipage} \\
\midrule

\textbf{6. 診斷外的價值} & 
\begin{minipage}[t]{9cm}
認識床邊診察超越診斷的力量\\
處理不確定性、建立關係、減少健康不平等
\end{minipage} \\

\end{longtable}

\section{策略一：前往床邊觀察}

\subsection{Sutton法則的應用}

\textcolor{blue}{\textbf{原始Sutton法則：}}「我搶銀行是因為錢在那裡」

\textcolor{blue}{\textbf{修訂版Sutton法則：}}要改善床邊臨床技能，就必須前往病人所在之處——床邊

\subsection{觀察技能的重要性}

\begin{longtable}{p{6cm}p{8cm}}
\caption{觀察在臨床診斷中的角色} \\
\toprule
\textbf{應用範疇} & \textbf{臨床價值} \\
\midrule
\endfirsthead
\textbf{應用範疇} & \textbf{臨床價值} \\
\midrule
\endhead
\bottomrule
\endfoot
\endlastfoot

診斷建立 & 
\begin{minipage}[t]{8cm}
- James Parkinson對震顫麻痺的描述幾乎完全依賴直接觀察\\
- 從床尾或走廊觀察即可發現關鍵診斷線索
\end{minipage} \\
\midrule

預後判斷 & 
\begin{minipage}[t]{8cm}
- 整體外觀反映疾病嚴重度\\
- 活動能力提供預後資訊
\end{minipage} \\
\midrule

個人狀況了解 & 
\begin{minipage}[t]{8cm}
- 觀察病人個人環境\\
- 理解社會支持系統\\
- 評估功能狀態
\end{minipage} \\
\midrule

教學評估 & 
\begin{minipage}[t]{8cm}
- 觀察學習者進行臨床診察\\
- 提供技能回饋的豐富機會\\
- 評估臨床能力發展
\end{minipage} \\

\end{longtable}

\subsection{提升觀察能力的方法}

\begin{enumerate}
    \item \textbf{臨床前期訓練：}
    \begin{itemize}
        \item 在非醫學情境練習觀察（例如：觀賞藝術）
        \item 可改善臨床觀察能力
        \item 建立系統化觀察習慣
    \end{itemize}
    
    \item \textbf{教師能力發展：}
    \begin{itemize}
        \item 透過練習提升直接觀察臨床技能的能力
        \item 增強教師信心
        \item 改善回饋品質
    \end{itemize}
\end{enumerate}

\section{策略二：實證理學檢查}

\subsection{假說驅動的理學檢查}

\begin{longtable}{p{4cm}p{10cm}}
\caption{實證理學檢查的步驟} \\
\toprule
\textbf{步驟} & \textbf{內容與方法} \\
\midrule
\endfirsthead
\textbf{步驟} & \textbf{內容與方法} \\
\midrule
\endhead
\bottomrule
\endfoot
\endlastfoot

\textbf{步驟1：估計檢查前機率} & 
\begin{minipage}[t]{10cm}
- 運用病史線索\\
- 考量疾病盛行率\\
- 建立診斷假說
\end{minipage} \\
\midrule

\textbf{步驟2：選擇檢查手法} & 
\begin{minipage}[t]{10cm}
- 針對疑似診斷選擇適當檢查\\
- 考慮檢查的診斷價值\\
- 如同選擇其他診斷工具般謹慎
\end{minipage} \\
\midrule

\textbf{步驟3：執行檢查} & 
\begin{minipage}[t]{10cm}
- 某些發現具病理診斷性（pathognomonic）\\
- 例如：帶狀疱疹的皮疹、蜂窩性組織炎的外觀與溫度
\end{minipage} \\
\midrule

\textbf{步驟4：與科技檢查比較} & 
\begin{minipage}[t]{10cm}
- 多數情況需與科技檢查結果比對\\
- 例如：心尖搏動外移與第三心音 vs. 心臟超音波\\
- 建立似然比（likelihood ratio）
\end{minipage} \\
\midrule

\textbf{步驟5：臨床決策} & 
\begin{minipage}[t]{10cm}
- 決定是否需要額外檢查\\
- 判斷資訊是否足以診斷\\
- 提供治療建議
\end{minipage} \\

\end{longtable}

\subsection{理學檢查的診斷價值}

\begin{itemize}
    \item \textcolor{blue}{\textbf{某些情況下，理學檢查仍是診斷標準}}
    \item \textcolor{blue}{\textbf{其他情況下，可利用似然比選擇適當檢查手法}}
    \item \textcolor{blue}{\textbf{透過與科技檢查比對，建立理學檢查的證據基礎}}
\end{itemize}

\section{策略三：創造刻意練習機會}

\subsection{刻意練習的重要性}

\begin{longtable}{p{5cm}p{9cm}}
\caption{刻意練習對臨床技能的影響} \\
\toprule
\textbf{面向} & \textbf{效果} \\
\midrule
\endfirsthead
\textbf{面向} & \textbf{效果} \\
\midrule
\endhead
\bottomrule
\endfoot
\endlastfoot

技能改善 & 
\begin{minipage}[t]{9cm}
- 專門臨床技能課程可改善理學檢查技能\\
- 結構化練習產生長期效果
\end{minipage} \\
\midrule

價值認知 & 
\begin{minipage}[t]{9cm}
- 與病人相處教導床邊技能的價值\\
- 資料收集立即影響病人照護\\
- 建立技能與臨床結果的連結
\end{minipage} \\
\midrule

克服障礙 & 
\begin{minipage}[t]{9cm}
- 賦予教師與學習者回到床邊的能力\\
- 打破「床邊沒有效率」的迷思
\end{minipage} \\

\end{longtable}

\subsection{晨間迴診的優化}

\begin{longtable}{p{4cm}p{10cm}}
\caption{以病人為中心的床邊迴診} \\
\toprule
\textbf{策略} & \textbf{實施方法} \\
\midrule
\endfirsthead
\textbf{策略} & \textbf{實施方法} \\
\midrule
\endhead
\bottomrule
\endfoot
\endlastfoot

事前準備 & 
\begin{minipage}[t]{10cm}
- 讓病人與學習者了解床邊迴診的期待\\
- 編排床邊互動流程\\
- 分配團隊成員任務
\end{minipage} \\
\midrule

角色分配 & 
\begin{minipage}[t]{10cm}
- 指定成員負責病歷輸入\\
- 指定成員負責資料檢索\\
- 確保每位成員有明確職責
\end{minipage} \\
\midrule

結構化呈現 & 
\begin{minipage}[t]{10cm}
- 聚焦臨床報告\\
- 保留時間澄清重要病史\\
- 將理學檢查細節整合進個案報告
\end{minipage} \\
\midrule

教學機會 & 
\begin{minipage}[t]{10cm}
- 運用結構化時間提供教學\\
- One-Minute Preceptor\\
- Five-Minute Moment
\end{minipage} \\
\midrule

效益 & 
\begin{minipage}[t]{10cm}
- 辨識立即需處理的問題\\
- 建立共同醫療議程\\
- 加速檢查、治療與會診的安排\\
- 提升醫師滿意度
\end{minipage} \\

\end{longtable}

\subsection{病人選擇策略}

\begin{enumerate}
    \item \textbf{高診斷價值情境：}
    \begin{itemize}
        \item 病史或理學檢查可能有助於診斷的情況
        \item 例如：慢性呼吸困難鑑別診斷（COPD vs. 心衰竭）
        \item 心肺疾病佔住院原因相當比例
    \end{itemize}
    
    \item \textbf{技能提升機會：}
    \begin{itemize}
        \item 選擇可改善教師與學習者臨床技能的個案
        \item 針對「神經恐懼症」(neurophobia)：選擇神經症狀病人
        \item 提供診斷與教學的雙重價值
    \end{itemize}
    
    \item \textbf{反應性學習：}
    \begin{itemize}
        \item 檢查臨床狀況新發生或惡化的病人
        \item 即時展示床邊診察在診斷與決策中的價值
        \item 工作場所非正式學習的重要部分
    \end{itemize}
\end{enumerate}

\subsection{專門教學時段}

\begin{longtable}{p{5cm}p{9cm}}
\caption{床邊技能教學時段設計} \\
\toprule
\textbf{時段類型} & \textbf{內容與方法} \\
\midrule
\endfirsthead
\textbf{時段類型} & \textbf{內容與方法} \\
\midrule
\endhead
\bottomrule
\endfoot
\endlastfoot

專門理學檢查時段 & 
\begin{minipage}[t]{9cm}
- 住院醫師提供欲探索的主訴或症狀\\
- 團隊至床邊進行理學檢查（包括超音波）\\
- 討論發現的臨床意義\\
- 與影像學等診斷結果比對
\end{minipage} \\
\midrule

晨會後床邊教學 & 
\begin{minipage}[t]{9cm}
- 晨會報告後至床邊\\
- 進行聚焦的病史詢問與理學檢查\\
- 床邊發現可能大幅改變鑑別診斷
\end{minipage} \\
\midrule

標準化病人教學 & 
\begin{minipage}[t]{9cm}
- 在傳統會議中使用標準化病人\\
- 示範與練習特定理學檢查手法\\
- 健康志願者協助
\end{minipage} \\
\midrule

教學效益 & 
\begin{minipage}[t]{9cm}
- 展示床邊臨床發現的價值\\
- 校正檢查技巧與影像結果\\
- 慶祝住院醫師的發現\\
- 承認床邊檢查的限制
\end{minipage} \\

\end{longtable}

\subsection{早期介入的重要性}

\begin{itemize}
    \item \textbf{習慣形成：}醫學生涯早期形成的習慣會影響整個職業生涯
    \item \textbf{臨床前教育：}
    \begin{itemize}
        \item 參與觀察性床邊診察（真實與標準化病人）
        \item 第三年實習結束時臨床技能較佳
    \end{itemize}
    \item \textbf{隱藏課程影響：}
    \begin{itemize}
        \item 臨床實習時未見醫師使用病史與理學檢查
        \item 學生快速轉向電子病歷與科技檢查
        \item 早期良好訓練可能被抵消
    \end{itemize}
\end{itemize}

\section{策略四：運用科技教學與強化臨床診察技能}

\subsection{數位聽診器}

\begin{longtable}{p{5cm}p{9cm}}
\caption{數位聽診器的功能與應用} \\
\toprule
\textbf{功能} & \textbf{臨床應用} \\
\midrule
\endfirsthead
\textbf{功能} & \textbf{臨床應用} \\
\midrule
\endhead
\bottomrule
\endfoot
\endlastfoot

歷史重要性 & 
\begin{minipage}[t]{9cm}
- Laënnec發明聽診器：兩世紀前最重要的醫學科技進展\\
- 現代科技擴展診斷能力
\end{minipage} \\
\midrule

即時視覺化 & 
\begin{minipage}[t]{9cm}
- 同時聽診與查看聲譜圖\\
- 即時顯示心電圖\\
- 增強教學效果
\end{minipage} \\
\midrule

人工智慧輔助 & 
\begin{minipage}[t]{9cm}
- AI演算法協助診斷心律不整\\
- 協助診斷瓣膜性心臟病\\
- 提高診斷準確性
\end{minipage} \\

\end{longtable}

\subsection{床邊超音波(POCUS)}

\begin{longtable}{p{5cm}p{9cm}}
\caption{POCUS在臨床教學中的角色} \\
\toprule
\textbf{面向} & \textbf{價值與應用} \\
\midrule
\endfirsthead
\textbf{面向} & \textbf{價值與應用} \\
\midrule
\endhead
\bottomrule
\endfoot
\endlastfoot

診斷能力 & 
\begin{minipage}[t]{9cm}
- 建立傳統理學檢查難以發現的診斷\\
- 即時影像化病理生理特徵
\end{minipage} \\
\midrule

技能校正 & 
\begin{minipage}[t]{9cm}
- 學習者校正理學檢查技能\\
- 連結檢查發現與即時影像\\
- 理解病理生理機轉
\end{minipage} \\
\midrule

AI輔助 & 
\begin{minipage}[t]{9cm}
- 協助影像擷取\\
- 協助影像判讀\\
- 提升學習曲線
\end{minipage} \\
\midrule

真正力量 & 
\begin{minipage}[t]{9cm}
- 將病人、醫師、學習者聚集在床邊\\
- 需要與病人互動\\
- 需要適當暴露檢查部位
\end{minipage} \\
\midrule

附加發現 & 
\begin{minipage}[t]{9cm}
- 過程中發現重要病史線索\\
- 發現理學檢查徵象（如胸骨切開疤痕、心律調節器）\\
- 促進醫病溝通
\end{minipage} \\
\midrule

共同決策 & 
\begin{minipage}[t]{9cm}
- 與病人建立連結\\
- 展示即時影像\\
- 促進參與式決策
\end{minipage} \\

\end{longtable}

\subsection{遠距醫療}

\begin{itemize}
    \item \textbf{克服障礙：}
    \begin{itemize}
        \item 克服個人防護裝備等物理障礙
        \item 限制觸診與聽診時提供即時診斷
        \item 病人臨床不穩定或感染管制需求時的替代方案
    \end{itemize}
    
    \item \textbf{重新定義床邊：}
    \begin{itemize}
        \item 包含遠距健康診察
        \item 病人自主進行醫師協助的理學檢查
        \item 提供醫師對病人居家環境的洞察
    \end{itemize}
\end{itemize}

\subsection{人工智慧的角色}

\begin{longtable}{p{5cm}p{9cm}}
\caption{AI在床邊診察的應用與限制} \\
\toprule
\textbf{應用領域} & \textbf{內容說明} \\
\midrule
\endfirsthead
\textbf{應用領域} & \textbf{內容說明} \\
\midrule
\endhead
\bottomrule
\endfoot
\endlastfoot

未來潛力 & 
\begin{minipage}[t]{9cm}
- 多模態AI系統協助觀察與視診\\
- 大型語言模型顯示臨床推理潛力
\end{minipage} \\
\midrule

立即影響 & 
\begin{minipage}[t]{9cm}
- 卸除電子病歷的行政負擔\\
- 釋放時間進行病人直接接觸\\
- 提升效率
\end{minipage} \\
\midrule

\textbf{關鍵原則} & 
\begin{minipage}[t]{9cm}
\textcolor{red}{\textbf{AI應告知而非取代：}}\\
- 人類觀察\\
- 人類臨床決策\\
- 人類溝通
\end{minipage} \\
\midrule

未來需求 & 
\begin{minipage}[t]{9cm}
- AI驅動的臨床工作場所更需要床邊臨床技能\\
- 強調直接觀察\\
- 強調技能回饋
\end{minipage} \\

\end{longtable}

\section{策略五：尋求並提供臨床技能回饋}

\subsection{床邊回饋的挑戰}

\begin{longtable}{p{4cm}p{10cm}}
\caption{床邊回饋的複雜性} \\
\toprule
\textbf{挑戰} & \textbf{考量因素} \\
\midrule
\endfirsthead
\textbf{挑戰} & \textbf{考量因素} \\
\midrule
\endhead
\bottomrule
\endfoot
\endlastfoot

情境複雜性 & 
\begin{minipage}[t]{10cm}
- 在病人面前提供回饋是複雜技能\\
- 可能涉及糾正理學檢查技巧錯誤\\
- 可能需要修改鑑別診斷或治療計畫
\end{minipage} \\
\midrule

潛在風險 & 
\begin{minipage}[t]{10cm}
- 回饋不當可能破壞學習者與病人關係\\
- 需要謹慎處理
\end{minipage} \\
\midrule

有效回饋 & 
\begin{minipage}[t]{10cm}
- 情境特定且周到的方式提供\\
- 可讓病人安心整個團隊投入照護\\
- 增強醫療團隊形象
\end{minipage} \\

\end{longtable}

\subsection{床邊回饋的最佳實務}

\begin{enumerate}
    \item \textbf{事前設定期待：}
    \begin{itemize}
        \item 提前告知病人與學習者將提供回饋
        \item 詢問學習者特定需求領域
        \item 建立心理準備
    \end{itemize}
    
    \item \textbf{避免醫學術語：}
    \begin{itemize}
        \item 盡可能使用淺顯易懂的語言
        \item 改善溝通效果
        \item 增進教育成效
    \end{itemize}
    
    \item \textbf{事後檢討：}
    \begin{itemize}
        \item 床邊診察後與學習者簡短檢討
        \item 強化關鍵要點
        \item 確認改善未來經驗的方法
    \end{itemize}
\end{enumerate}

\subsection{直接觀察的現況}

\begin{longtable}{p{5cm}p{9cm}}
\caption{美國與其他國家的臨床技能評估差異} \\
\toprule
\textbf{地區/方式} & \textbf{評估方法與效果} \\
\midrule
\endfirsthead
\textbf{地區/方式} & \textbf{評估方法與效果} \\
\midrule
\endhead
\bottomrule
\endfoot
\endlastfoot

\textbf{美國現況} & 
\begin{minipage}[t]{9cm}
- 真實病人的直接觀察與臨床技能評估罕見\\
- 依賴多選題考試評估醫學知識\\
- 判斷學習者發展中的能力
\end{minipage} \\
\midrule

\textbf{其他國家} & 
\begin{minipage}[t]{9cm}
- 內科住院醫師畢業需通過實際病人考試\\
- 真實病人呈現真實發現\\
- 教師觀察員評估
\end{minipage} \\
\midrule

\textbf{效果差異} & 
\begin{minipage}[t]{9cm}
- 高風險評估可驅動學習\\
- 非美國訓練醫師臨床檢查技能通常更熟練\\
- 對床邊診察價值有更廣泛認識\\
- 可能轉化為更佳病人結果
\end{minipage} \\

\end{longtable}

\subsection{APECS評估系統}

\begin{longtable}{p{4cm}p{10cm}}
\caption{理學檢查與溝通技能評估(APECS)} \\
\toprule
\textbf{特徵} & \textbf{內容} \\
\midrule
\endfirsthead
\textbf{特徵} & \textbf{內容} \\
\midrule
\endhead
\bottomrule
\endfoot
\endlastfoot

評估性質 & 
\begin{minipage}[t]{10cm}
- 內科住院醫師形成性評估經驗\\
- 包含真實病人的整合式診察\\
- 僅理學檢查診察\\
- 遠距醫療診察\\
- POCUS診察
\end{minipage} \\
\midrule

評估領域 & 
\begin{minipage}[t]{10cm}
- 跨七個臨床領域評估住院醫師\\
- 經驗豐富的教師提供個別化實作回饋
\end{minipage} \\
\midrule

\textbf{重要發現} & 
\begin{minipage}[t]{10cm}
\textcolor{blue}{\textbf{技巧與發現的關聯：}}\\
- 良好技巧與正確識別理學檢查發現顯著相關\\
- 良好技巧與正確鑑別診斷顯著相關\\
- 正確識別發現與正確鑑別診斷顯著相關
\end{minipage} \\
\midrule

實際案例 & 
\begin{minipage}[t]{10cm}
- 嚴重主動脈瓣逆流病人檢查\\
- 僅半數實習醫師聽到特徵性舒張期雜音\\
- 常見錯誤：隔著衣服聽診\\
- 常見錯誤：聽診時摸橈動脈而非心尖搏動或頸動脈
\end{minipage} \\
\midrule

教師發展 & 
\begin{minipage}[t]{10cm}
- 不同專科與經驗的醫師配對為指導教師\\
- 與學習者共同學習\\
- 提供教師發展機會
\end{minipage} \\

\end{longtable}

\subsection{實務工具}

\begin{longtable}{p{5cm}p{9cm}}
\caption{臨床技能觀察與回饋工具} \\
\toprule
\textbf{工具} & \textbf{應用方式} \\
\midrule
\endfirsthead
\textbf{工具} & \textbf{應用方式} \\
\midrule
\endhead
\bottomrule
\endfoot
\endlastfoot

Ten-Minute Moment & 
\begin{minipage}[t]{9cm}
- 床邊醫學學會開發\\
- 日常工作流程中使用\\
- 教師觀察學習者進行聚焦理學檢查\\
- 使用技能工作表提供即時回饋
\end{minipage} \\
\midrule

Mini-CEX & 
\begin{minipage}[t]{9cm}
- Mini-Clinical Evaluation Exercise\\
- 常用於指導工作場所觀察\\
- 結構化回饋工具\\
- 廣泛使用
\end{minipage} \\

\end{longtable}

\section{策略六：認識床邊診察超越診斷的力量}

\subsection{處理不確定性}

\begin{longtable}{p{5cm}p{9cm}}
\caption{臨床不確定性的處理} \\
\toprule
\textbf{面向} & \textbf{策略與價值} \\
\midrule
\endfirsthead
\textbf{面向} & \textbf{策略與價值} \\
\midrule
\endhead
\bottomrule
\endfoot
\endlastfoot

不確定性本質 & 
\begin{minipage}[t]{9cm}
- 臨床醫學在不確定性背景下實踐\\
- 存在於診斷、治療決策、臨床追蹤各階段\\
- 是床邊病人照護的一部分
\end{minipage} \\
\midrule

承認的價值 & 
\begin{minipage}[t]{9cm}
- 承認不確定性的存在\\
- 允許病人、臨床醫師、學習者共同協商\\
- 可強化醫病關係
\end{minipage} \\
\midrule

好奇心的角色 & 
\begin{minipage}[t]{9cm}
- Fitzgerald (1999)：「好奇心是調查、發現的衝動」\\
- 應用於臨床醫學：將病人視為一個人\\
- 徵象與症狀是故事中的片段\\
- 減緩不確定性的重要屬性
\end{minipage} \\
\midrule

夥伴關係 & 
\begin{minipage}[t]{9cm}
- 以好奇心接近床邊診察\\
- 有效與病人和學習者建立夥伴關係\\
- 共同面對臨床挑戰
\end{minipage} \\

\end{longtable}

\subsection{完全在場的實證做法}

\begin{longtable}{p{4cm}p{10cm}}
\caption{與病人互動時完全在場的策略} \\
\toprule
\textbf{做法} & \textbf{效果} \\
\midrule
\endfirsthead
\textbf{做法} & \textbf{效果} \\
\midrule
\endhead
\bottomrule
\endfoot
\endlastfoot

有意圖的準備 & 
\begin{minipage}[t]{10cm}
- Preparing with intention\\
- 心理與實體準備\\
- 建立正念狀態
\end{minipage} \\
\midrule

專注傾聽 & 
\begin{minipage}[t]{10cm}
- Listening intently\\
- 全神貫注於病人敘述\\
- 避免分心
\end{minipage} \\
\midrule

共識重點 & 
\begin{minipage}[t]{10cm}
- Agreeing on what matters\\
- 與病人建立共同目標\\
- 確立優先事項
\end{minipage} \\
\midrule

\textbf{效益} & 
\begin{minipage}[t]{10cm}
\textcolor{blue}{\textbf{超越臨床線索的價值：}}\\
- 與病人故事連結\\
- 建立信任\\
- 尋找工作意義\\
- 對抗壓力與職業倦怠的有力支撐
\end{minipage} \\

\end{longtable}

\subsection{理學檢查的療癒效果}

\begin{itemize}
    \item \textcolor{blue}{\textbf{傳達關懷：}}適當執行的理學檢查可傳達關懷
    \item \textcolor{blue}{\textbf{安慰劑效應：}}超越診斷發現的安慰劑效果
    \item \textcolor{blue}{\textbf{慶祝發現：}}
    \begin{itemize}
        \item 慶祝床邊發現產生重要臨床資訊的時刻
        \item 在臨床診察中建立興奮感
        \item 對抗「靈光乍現缺乏症」(eurekapenia)
    \end{itemize}
    \item \textbf{Eurekapenia定義：}
    \begin{itemize}
        \item 缺乏「靈光乍現」時刻
        \item 學習者認識臨床發現與疾病病理生理基礎之間連結
        \item 許多學習者與執業醫師經歷此現象
        \item 原因：很少時間在床邊
    \end{itemize}
\end{itemize}

\subsection{解決健康照護差異}

\begin{longtable}{p{5cm}p{9cm}}
\caption{床邊診察在減少健康不平等的角色} \\
\toprule
\textbf{問題} & \textbf{床邊診察的貢獻} \\
\midrule
\endfirsthead
\textbf{問題} & \textbf{床邊診察的貢獻} \\
\midrule
\endhead
\bottomrule
\endfoot
\endlastfoot

種族與族裔差異 & 
\begin{minipage}[t]{9cm}
- 少數族裔青少年從未接受常規體檢的比例高於白人同儕\\
- 即使檢查，系統性不平等可能惡化結果
\end{minipage} \\
\midrule

教育資源不足 & 
\begin{minipage}[t]{9cm}
- 包含深色膚色的皮膚科教科書缺乏\\
- 脈搏血氧儀等裝置存在偏差\\
- 可能導致有色人種病人結果惡化
\end{minipage} \\
\midrule

語言障礙 & 
\begin{minipage}[t]{9cm}
- 英語能力有限的病人可能接受不適當照護\\
- 溝通障礙影響診斷準確性
\end{minipage} \\
\midrule

\textbf{教育者的角色} & 
\begin{minipage}[t]{9cm}
\textcolor{blue}{\textbf{帶團隊至床邊的機會：}}\\
- 以情境特定方式承認這些不平等\\
- 解決健康照護差異\\
- 提供教學示範
\end{minipage} \\

\end{longtable}

\section{核心要點總結}

\begin{longtable}{p{14cm}}
\caption{重振床邊臨床診察的關鍵訊息} \\
\toprule
\textbf{核心要點} \\
\midrule
\endfirsthead
\textbf{核心要點} \\
\midrule
\endhead
\bottomrule
\endfoot
\endlastfoot

\textcolor{red}{\textbf{1. 現況問題：}}21世紀醫學學習者訓練期間與病人相處時間少於20世紀同儕，導致床邊臨床技能知識與實踐下降 \\
\midrule

\textcolor{red}{\textbf{2. 嚴重後果：}}床邊臨床技能下降導致診斷錯誤、臨床結果不良、醫療成本增加、醫師職業倦怠 \\
\midrule

\textcolor{red}{\textbf{3. 觀察與練習：}}帶學習者至床邊促進臨床觀察技能、創造練習機會、允許實證示範檢查技能 \\
\midrule

\textcolor{red}{\textbf{4. 科技整合：}}床邊診察整合床邊科技與人工智慧，補充人類觀察、人類臨床決策、人類溝通 \\
\midrule

\textcolor{red}{\textbf{5. 情境回饋：}}尋求機會以情境特定方式提供臨床技能回饋，改善床邊技術與從診察獲得資訊的詮釋 \\
\midrule

\textcolor{red}{\textbf{6. 多重價值：}}超越床邊臨床診察獲得的診斷資料，理學檢查幫助學習者處理臨床不確定性、幫助教師示範與病人互動、改善醫病溝通、增加專業成就感、協助解決健康照護差異 \\

\end{longtable}

\section{結論}

在科技進步、與病人時間有限、臨床不確定性的背景下，迫切需要重振床邊診察以滿足21世紀病人、學習者與臨床教育者的需求。透過運用六大策略，臨床教育者可以幫助學習者認識床邊診察在診斷推理中的價值、強化醫病關係、對抗健康照護不平等、改善專業成就感並避免職業倦怠。

Osler一個多世紀前的話語至今仍然真實：「醫學是在床邊學習，而非在教室中學習。不要讓你對疾病表現的概念來自講堂聽到的話語或從書本讀到的文字。觀察，然後推理、比較與控制。但首先要觀察。」

\section{參考文獻}

\begin{enumerate}[leftmargin=*]

\item Rosen MA, Bertram AK, Tung M, Desai SV, Garibaldi BT. Use of a real-time locating system to assess internal medicine resident location and movement in the hospital. \href{https://doi.org/10.1001/jamanetworkopen.2022.15885}{\textit{JAMA Netw Open}. 2022;5(6):e2215885.}

%\item Vukanovic-Criley JM, Hovanesyan A, Criley SR, et al. Confidential testing of cardiac examination competency in cardiology and noncardiology faculty and trainees: a multicenter study. \href{https://doi.org/10.1002/clc.20843}{\textit{Clin Cardiol}. 2010;33:738-45.}

\item Newman-Toker DE, McDonald KM, Meltzer DO. How much diagnostic safety can we afford, and how should we decide? A health economics perspective. \href{https://doi.org/10.1136/bmjqs-2012-001616}{\textit{BMJ Qual Saf}. 2013;22(Suppl 2):ii11-ii20.}

\item Singh H, Giardina TD, Meyer AND, Forjuoh SN, Reis MD, Thomas EJ. Types and origins of diagnostic errors in primary care settings. \href{https://doi.org/10.1001/jamainternmed.2013.2777}{\textit{JAMA Intern Med}. 2013;173:418-25.}

\item Garibaldi BT, Elder A. Seven reasons why the physical examination remains important. \href{https://doi.org/10.4997/jrcpe.2021.301}{\textit{J R Coll Physicians Edinb}. 2021;51:211-4.}

\item Neumann M, Edelhäuser F, Tauschel D, et al. Empathy decline and its reasons: a systematic review of studies with medical students and residents. \href{https://doi.org/10.1097/ACM.0b013e318221e615}{\textit{Acad Med}. 2011;86:996-1009.}

%\item Costanzo C, Verghese A. The physical examination as ritual: social sciences and embodiment in the context of the physical examination. \href{https://doi.org/10.1016/j.mcna.2018.01.003}{\textit{Med Clin North Am}. 2018;102:425-31.}

%\item Cooper LA, Roter DL. Patient-provider communication: the effect of race and ethnicity on process and outcomes of healthcare. In: Smedley BD, Stith AY, Nelson AR, eds. \textit{Unequal treatment: confronting racial and ethnic disparities in health care}. Washington, DC: National Academies Press, 2004:552-93.

%\item Wu EH, Fagan MJ, Reinert SE, Diaz JA. Self-confidence in and perceived utility of the physical examination: a comparison of medical students, residents, and faculty internists. \href{https://doi.org/10.1007/s11606-007-0387-0}{\textit{J Gen Intern Med}. 2007;22:1725-30.}

%\item Barry JM. \textit{The great influenza: the story of the deadliest pandemic in history}. New York: Penguin Books, 2004.

\item McGee S. \textit{Evidence-based physical diagnosis}. 5th ed. Philadelphia: Elsevier, 2022.

\item Verghese A. Culture shock — patient as icon, icon as patient. \href{https://doi.org/10.1056/NEJMp0807461}{\textit{N Engl J Med}. 2008;359:2748-51.}

%\item Block L, Habicht R, Wu AW, et al. In the wake of the 2003 and 2011 duty hours regulations, how do internal medicine interns spend their time? \href{https://doi.org/10.1007/s11606-013-2453-3}{\textit{J Gen Intern Med}. 2013;28:1042-7.}

\item Gonzalo JD, Chuang CH, Huang G, Smith C. The return of bedside rounds: an educational intervention. \href{https://doi.org/10.1007/s11606-010-1344-7}{\textit{J Gen Intern Med}. 2010;25:792-8.}

%\item Garibaldi BT, Russell SW. Strategies to improve bedside clinical skills teaching. \href{https://doi.org/10.1016/j.chest.2021.04.078}{\textit{Chest}. 2021;160:2187-95.}

\item Reilly BM. Physical examination in the care of medical inpatients: an observational study. \href{https://doi.org/10.1016/S0140-6736(03)14464-9}{\textit{Lancet}. 2003;362:1100-5.}

\item Verghese A, Charlton B, Kassirer JP, Ramsey M, Ioannidis JPA. Inadequacies of physical examination as a cause of medical errors and adverse events: a collection of vignettes. \href{https://doi.org/10.1016/j.amjmed.2015.06.004}{\textit{Am J Med}. 2015;128(12):1322-1324.e3.}

%\item Garibaldi BT, Niessen T, Gelber AC, et al. A novel bedside cardiopulmonary physical diagnosis curriculum for internal medicine postgraduate training. \href{https://doi.org/10.1186/s12909-017-1022-7}{\textit{BMC Med Educ}. 2017;17:182.}

\item Holmboe ES, Hawkins RE, Huot SJ. Effects of training in direct observation of medical residents' clinical competence: a randomized trial. \href{https://doi.org/10.7326/0003-4819-140-11-200406010-00008}{\textit{Ann Intern Med}. 2004;140:874-81.}

%\item Elder A. Clinical skills assessment in the twenty-first century. \href{https://doi.org/10.1016/j.mcna.2018.03.001}{\textit{Med Clin North Am}. 2018;102:545-58.}

\end{enumerate}

\vspace{1cm}
\textbf{關鍵詞：}床邊醫學、臨床技能、理學檢查、醫學教育、病人照護、直接觀察

\end{document}